% 04-walls.tex: the schema as the security boundary.
\section{The walls: the schema as the security boundary}
\label{sec:walls}

\subsection{Why the database is not plumbing}

In most systems the database stores what the application decided. In Polaris the database decides,
and the application is one client among several: the operator's web application, the command-line
tool, a bulk loader, a database shell, a restore from backup. A rule that lives in application code
binds only the clients that run that code. A rule enforced by a trigger, a check constraint or a
unique index binds every client, survives every restore, and cannot be bypassed by the next caller,
because the application role has no data-definition privilege with which to remove it. The
governing sentence of the architecture is \emph{by construction}: the system cannot be operated in a
way that violates the constraint, because the constraint is enforced before any application code
runs. \sImpl

The design principle behind it is older than Polaris: make invalid states impossible to represent
rather than asking developers not to create them. A verification event that records a token
identifier at the zero-knowledge disclosure level is not an application bug to be caught in review.
It is a privacy violation that must never be storable, so it is a check constraint on the row.

\subsection{The ten constraints and where each lives}

The constitution names ten hard constraints in two tiers. Five are constitutional, rights
guarantees that may not be relaxed without an amendment recorded in the changelog; five are
engineering invariants, the disciplines that keep the first five true under load and attack. The
hierarchy is one line deep, vocation, then constitutional constraints, then engineering invariants,
and the constitution states that no deeper apparatus governs them.

\begin{table}[H]
\centering\footnotesize
\caption{The ten constraints at \code{v9.345}, their tier, and the object that enforces each. The table is kept honest against the code by \code{check\_c1c10\_objects\_resolve}, which fails the build if a named object no longer exists.}
\label{tab:constraints}
\setlength{\tabcolsep}{4pt}\renewcommand{\arraystretch}{1.1}
\rowcolors{2}{tabstripe}{white}
\begin{tabularx}{\textwidth}{>{\bfseries}p{0.3in} >{\RaggedRight\arraybackslash}X >{\RaggedRight\arraybackslash}p{0.95in} >{\RaggedRight\arraybackslash}p{2.0in}}
\toprule
\rowcolor{tabheader}\thead{\#} & \thead{Constraint} & \thead{Tier} & \thead{Enforced by} \\
\midrule
C1 & The audit trail is append-only; history is never rewritten or deleted & Constitutional & Triggers reject UPDATE and DELETE on every audit table (\code{reject\_audit\_modification}) \\
C2 & A zero-knowledge verification stores no token identifier & Constitutional & Bidirectional CHECK \code{chk\_disclosure\_token\_consistency} \\
C3 & One person holds at most one ACTIVE token & Constitutional & Partial unique index \code{uq\_one\_active\_per\_person} \\
C4 & Failed-login counting is atomic & Engineering & Single-statement \code{UPDATE \ldots RETURNING} \\
C5 & No inline scripts; the content security policy is \code{script-src 'self'} & Engineering & Response header, verified per route \\
C6 & Disclosure level is enforced server-side; a client cannot upgrade what it learns & Constitutional & Server code paired with redaction tests; the C2 CHECK behind it \\
C7 & No hardcoded cryptography; algorithms are rows in a registry & Engineering & Foreign key to \entity{CryptographicAlgorithm} \\
C8 & Every map and API aggregate is bounded & Engineering & Hard caps in the SQL functions and routes \\
C9 & Concurrency claims are tested with real threads & Engineering & Threaded suites against a live database \\
C10 & Identity is not money; the schema carries no monetary claim & Constitutional & Structural absence, pinned by a check \\
\bottomrule
\end{tabularx}
\end{table}

The repository also records a structural claim about the ten: that the set is closed, so that an
eleventh constraint requires either replacing one or extending the framework, and that the
constraints depend on one another, so that removing any one ripples through the rest. The
dependency walk (remove C10 and every other constraint protects a financial-surveillance system
instead; remove C1 and every other constraint's enforcement becomes retroactively deniable) is the
argument for treating them as a graph rather than a list. Athena, in Section~\ref{sec:cognition},
holds the map from each constraint to the exact live mechanism that enforces it, and a check fails
the build if a mechanism drifts.

\figfile{districts}{The 36 canonical tables of \code{01\_schema.sql} at \code{v9.345}, as rings
around the Version 1 nucleus. The faint lattice at the centre is the nucleus geometry of
Figure~\ref{fig:nucleus} at the same proportions. Ring two is the credential machinery: signatures,
permissions, epochs, recovery, duress. Ring three is authority and trust: who may sign under which
algorithm, who trusts whom, revocation ceilings, quotas, relying parties, the replay registers. Ring
four is the archives and watchtowers: append-only logs, anchors, the retention decision as data. A
migrated deployment holds seven more tables (the migration registry, operator sessions and
credentials, the audit-access log, and Athena's three curated tables), and four partition
children.}{fig:districts}

\subsection{The audit of record}

Thirteen tables follow one named pattern, the audit of record: the row and its history are the same
object. There is no separate event log beside the table that could be edited while the table is
locked. Mutation is either forbidden outright or bounded to one field that moves one way, such as a
signature's deprecation date or an attestation's revocation date. Every one of the thirteen is
enforced by a trigger that raises \code{insufficient\_privilege}, and a check fails the build if any
trigger stops doing so. \entity{RecoveryRequest} was the one instance resting on a partial unique
index and procedure discipline, so a raw update from a database session was not refused; at
\code{v9.347} its decision fields moved one way under a trigger like the other twelve, and the check
now covers thirteen of thirteen. \sImpl

The corollary is that no foreign key anywhere in the schema cascades a delete. A cascade would let
the lifecycle events of a deleted token vanish with it. Every reference defaults to \code{NO
ACTION}, so a principal referenced by any audit row is effectively undeletable, and a check scans
every SQL file for both cascade forms with no allowlist.

The application role, \code{polaris\_app}, is revoked UPDATE and DELETE on the append-only tables
and holds no data-definition privilege. The only path that deletes audit rows is the archive-then-purge
procedure, which verifies the archive against its manifest first, refuses any cutoff inside the
retention window, and appends its own append-only checkpoint row. The retention window itself is
data with a floor: no policy row can express a retention shorter than a year, because a check
constraint refuses it, and shortening the floor is a schema change answered for in review. The
vocation refuses unbounded retention; it equally refuses retention short enough to erase the record.
\sImpl

\subsection{The check layer, and the checks that check the checks}

Above the schema sits a flat layer of 191 plain functions, each of the form
\code{check\_*(repo\_root)}, that read the repository and return a finding. They gate continuous
integration: any failure fails the build. Each check is paired with a detection test that breaks a
fixture in the specific way the check claims to catch and asserts that the check turns red. A check
that cannot detect its own violation is treated as broken. This is the mechanism by which prose
claims are turned into gates: a count stated in the README, a route that must return a header, a
migration that must not cascade, a drill that CI must run. \sImpl

\begin{layercard}{Layer card: the walls}
\qa{What it is}{PostgreSQL 16: 36 tables, 26 triggers, the stored procedures that implement every state-changing use case, the append-only audit tables, the retention policy as data.}
\qa{Why it exists}{A promise enforced only by the people who run the system is not a promise. The schema binds every client, every restore and every insider with a database connection short of the owner.}
\qa{What it trusts}{The database engine and its owner; the correctness of the constraints as written.}
\qa{What it refuses}{Application code, review process and operator discipline as the last line: the application never enforces a constitutional guarantee on its own.}
\qa{What it can see}{Everything the system stores. This is why the privacy layer bounds what is stored, not who may read it.}
\qa{What it cannot see}{Biometric templates, holder-side keys, the content of exchanged messages or timestamped documents: none enter the database.}
\qa{If it fails}{A dropped trigger, a lowered retention floor or a cascade fails the build through the checks; a database owner with a shell has larger options than any trigger, which the threat model states rather than hides.}
\qa{Checked by}{\code{check\_aor\_append\_only\_triggers}, \code{check\_aor\_privilege\_boundary}, \code{check\_no\_fk\_cascade}, \code{check\_retention\_engine}, the SQL self-tests, the CHECK-constraint suite.}
\qa{Now or later}{Now.}
\qa{Inside and outside}{Inside: the credential core. Outside: cryptographic evidence, because a constraint binds only the database that holds it.}
\end{layercard}

\figfile{town-map}{The walled town, with the real component beside every civic name. The register
in the hall is the credential core; the wall is the schema; the four gates are the routes through
which a claim enters or leaves; the watchtowers are the transparency logs and the witnesses outside
the wall watch them; the road to the second town is a directional, per-context attestation; the laws
above the town bind the town itself. The plan describes structure. Polaris is federated: every
authority is its own town.}{fig:town-map}

\nextlayer{A constraint binds the database, not a copy of the claim that has left it. The moment a
credential is handed to a stranger, or an authority's history is read by an auditor, nothing in the
schema travels with it. The next layer makes claims carry their own proof.}
